\documentclass[../main.tex]{subfiles}
\begin{document}
%==========================================TIÊU ĐỀ TẬP========================================
\begin{table}[h]
  \begin{adjustwidth}{-1cm}{}
    \begin{tabular}{>{\centering\arraybackslash}p{6cm}|>{\raggedright\arraybackslash}p{12.5cm}}
      \multirow{5}{6cm}
      % Cột bên trái: ÉPISODE và số tập
      {\\[-40pt]
        \flaregothic\fontsize{20pt}{20pt}\selectfont \centering
        {\contour{errorfive}{\textcolor{black}{ÉPISODE}}}\color{black}\\
        \burbank\fontsize{72pt}{72pt}\selectfont
        \includegraphics[height=3.12359cm]{../../../../Image/Book?/number/num29.png}
        \\[18pt]
        % \flaregothic\fontsize{14pt}{14pt}\selectfont
        % \color{epattr}BIENVENUE, \uline{\color{Banpire} Mel} !
        % \flaregothic\fontsize{18pt}{18pt}\selectfont
        % \color{epattr}ÉPISODE PERDU
        \color{black}
      }
       &
      % Dòng thứ nhất: tên tập tiếng Nhật
      {\fotskip\fontsize{18pt}{18pt}\selectfont \ruby{幽}{ゆう}\ruby{霊}{れい}の\ruby{森}{もり}の\ruby{冒}{ぼう}\ruby{険}{けん}}
      \\[6pt]
      % Dòng thứ hai: tên tập tiếng Anh
       & {\cafeteria\fontsize{18pt}{18pt}\selectfont Haunted Woods Adventure}                                                                                       \\[4pt]
      % Dòng thứ ba: tên tập tiếng Việt
       & {\vnmsans\fontsize{18pt}{18pt}\selectfont Phiêu lưu trong rừng ma ám}                                                                                      \\[8pt]
      % Dòng thứ tư: nhân vật xuất hiện
       & {\caviar{\fontsize{14pt}{14pt}\selectfont Nhân vật: \emp{azki} \emp{miko} \emp{aki} \emp{mel} \emp{aqua} \emp{shion} \emp{marine} \emp{lamy} \emp{polka}}}
      % \\[6pt]
      % Dòng cuối cùng: Thời gian diễn ra của tập (thường sẽ là ngày phát hành)
      % & {\continuum\fontsize{16pt}{16pt}\selectfont Ngày phát hành: 11/06/2022}\\[8pt]
    \end{tabular}
  \end{adjustwidth}
\end{table}
%===========================================NỘI DUNG==========================================
\color{black}

\begin{center}
  \pov{Hikawa Suzuki}
  \uline{Ngày hôm sau - Chín giờ tối.}
\end{center}

Tuyến đường về phía khu rừng cổ tối tăm hơn tôi hình dung.

Ánh đèn đường thưa dần khi chúng tôi xa khu dân cư, chỉ còn vệt vàng nhạt trên mặt nhựa ướt lạnh. Phía trước, rừng già đen kịt trải xuống sườn đồi như tấm màn bí ẩn.

Cả hội tách thành từng nhóm nhỏ.

Azuki-san bước phía trước với vẻ hứng thú rõ ràng, đôi mắt sáng lên như thể sắp bước lên sân khấu. Tôi đoán cô ấy đang tưởng tượng đến thứ âm thanh bí ẩn mà mình sắp được nghe.

Ngay sau cô, Honoka-san ríu rít khoe chuyện, nói đủ trên trời dưới biển, còn Shion-san khúc khích như thể cả chuyến đi này không khác nào một thú vui nghịch ngợm.

Aku-san thì đi cạnh họ với vẻ hơi miễn cưỡng, hai tay khoanh trước ngực. ``Nhắc lại lần nữa,'' cậu ấy lầm bầm, ``m-mình chỉ đi vì mấy cậu kéo mình theo thôi đấy.''

``Biết rồi, nhưng ngươi vẫn đi đấy thôi,'' Shion-san cười đáp. ``Kể cả ngươi không muốn thì ta vẫn sẽ lôi ngươi theo thôi, thuộc hạ của ta.''

Phía sau họ một chút là Touka-san, Saki-san, Mari-san và Saya-san. Bốn người họ đang đi với mỗi tâm trạng khác nhau.

Cô bạn tóc xanh quan sát tỉ mỉ như ghi số liệu thí nghiệm. Cô bạn tóc vàng thỉnh thoảng cúi xem bụi cỏ ven đường.

Cô bạn tóc đỏ ít nói, bước chậm, nét mặt có gì đó suy tư. Cô bạn tóc ngắn thì ngẩng sao trời, như đang đọc bản đồ thiên hà.

Cuối đoàn là bốn chúng tôi: Onii-sama, Onee-sama, Sakura-san và tôi.

Tiếng côn trùng rỉ rả từ rừng. Gió lạnh cuối xuân đêm xuống khiến chúng tôi khẽ run rẩy.

Kỳ lạ thật.

Chỉ ba ngày trước thôi, tôi vẫn chưa từng nghĩ mình sẽ đi bộ tới một khu rừng lúc chín giờ tối để tìm tiếng hát của những đứa trẻ vô hình.

Tất cả mọi chuyện bắt nguồn từ lời đồn trong lớp\dots

\begin{center}
  \uline{Hồi tưởng.\\Ba ngày trước - giờ nghỉ giữa giờ.}
\end{center}

Buổi học thứ hai, môn Vật lý, vừa kết thúc. Tiếng chuông ngân dài, cửa lớp mở toang, cả lũ như cá đội nước tràn ra hành lang. Lớp 1-C, ô cửa cuối khẽ hé, gió tháng Sáu thổi vào mùi cỏ non và ẩm đất bãi trống sau trường. Nắng cao, nhưng kính chia thành ô vuông xám trắng, đè lên mặt bàn gỗ lấm bấm vết chì và tẩy.

Cả phòng học ồn ào như thường lệ khi giờ nghỉ bắt đầu: tiếng bật nắp bình nước, tiếng ghế kéo lê sàn, tiếng ai đó đập cặp xuống bàn để đuổi con ruồi vo ve quanh cốc sữa chua dâu. Tôi đang định lấy sách ra khỏi cặp để đọc thì nghe thấy tiếng Honoka-san nói lớn ở phía bàn gần cửa sổ.

``\textbf{Này, này! Có tin hot đây!}''

Câu nói của cô ấy luôn có sức hút kỳ lạ. Chỉ cần Honoka-san lên tiếng, vài người xung quanh đã tự nhiên quay lại. Đôi lúc cũng không phải tốt lành gì, vì cô ấy toàn bị mang danh là ``nàng hề'' của lớp khi toàn nói những thứ vô nghĩa hoặc đề xuất một ý tưởng gì đó ngớ ngẩn.

``Gì thế?'' Shion-san hỏi.

Honoka-san hạ giọng xuống một chút, nhưng vẻ mặt thì rõ ràng đang rất thích thú. ``Đàn anh lớp trên của tớ kể rằng tối hôm kia có người nghe thấy tiếng trẻ con hát trong rừng phía sau đền cũ đấy!''

Một vài người trong lớp bật cười.

``Chắc là mấy đứa tiểu học chơi trốn tìm thôi,'' Nanase-san lên tiếng. ``Không đáng quan tâm.''

Nhưng cô Phù thủy nghiêng đầu: ``Tiếng hát ư?''

``Ừ,'' cô bạn tai Cáo sa mạc gật đầu, ``nhưng vấn đề là\dots\ lúc đó đã hơn chín giờ tối rồi. Phát ra ở trong rừng, và không có ai ở đó vào giờ muộn đó cả.''

Lớp học bỗng chốc trùng xuống. Những tiếng xì xào bắt đầu vang lên. Có người thì cho rằng đây chỉ là trò đùa nhảm của học sinh cấp dưới, có người lại nghĩ đến chuyện ma quái rợn người, nhưng tất cả đều không giấu được sự tò mò trước lời đồn đáng sợ ấy.

Azuki-san chợt lên tiếng: ``\dots Tiếng hát sao?''

Cô ấy quay hẳn người lại từ bàn phía trước. Ánh mắt của cô ấy lúc đó\dots\ giống hệt khi cô ấy nói về âm nhạc trong giờ sinh hoạt câu lạc bộ. ``Tớ muốn nghe thử.''

Ngay khi cô ấy dứt lời, mọi ánh nhìn đều đổ dồn về phía cô ấy.

``Cậu nói gì cơ?'' Honoka-san bật cười.

Azuki-san nhún vai: ``Nếu có người hát, tớ muốn biết chúng hát bài gì.''

Shion-san giơ ngón tay cái lên, cười gian: ``Nghe có vẻ thú vị đấy!''

Tôi thấy Onii-sama thở dài bên bàn. ``Azuki-san,'' anh nói, giọng không giấu được sự mỉa mai, ``cậu có vẻ thích thú với chuyện này lắm nhỉ?''

``Thì sao?'' Azuki-san cười.

Shion-san lập tức chen vào, nhoẻn miệng: ``Không phải mấy ngươi sợ đấy chứ? Thuộc hạ của ta, ngươi nghĩ sao?''

Aku-san nhăn mặt. ``Đừng có kéo mình vào mấy trò này.''

``Nhưng mà ngươi vẫn đang nghe mà,'' Shion-san đáp.

Ở phía sau, Touka-san đi tới. ``Nếu thật sự có người nghe thấy âm thanh trong rừng thì việc kiểm chứng cũng không phải là ý tồi.''

Saki-san cũng quay lại, ánh mắt tỏ lên sự hứng thú nhẹ: ``Mình nghe nói trong khu rừng đó có vài loại hoa dại hiếm. Nếu đi vào ban đêm, có thể sẽ thấy chúng nở.''

Mari-san gật đầu nhẹ: ``Mình cũng muốn xem thử phong cảnh trong rừng lúc đó. Ánh sáng ban đêm đôi khi tạo ra những bố cục rất thú vị.''

Tôi không biết tại sao, nhưng khi cô nàng Mỹ thuật nói câu đó, tôi có cảm giác cô ấy đang giấu một điều gì khác. Một cái gì đó mà cậu ấy đã từng trải qua trước đây.

Cuối cùng, Saya-san khẽ lên tiếng: ``Khu rừng phía sau đền cũ\dots\ nếu mình nhớ đúng thì ngày xưa đó là vị trí của một đàn tế cổ.''

Cả lớp lập tức sôi nổi lên.

``Khoan đã, nghe đáng sợ rồi đấy!''

``Đàn tế cơ à?''

Honoka-san thì lại càng phấn khích hơn, như một đứa trẻ tìm thấy bản đồ kho báu:``Quá hợp luôn! Tớ muốn đi bây giờ!''

Onii-sama chống cằm nhìn cảnh đó một lúc. Rồi anh khẽ nói với tôi: ``\vie{Cách Azuki-san tập hợp mọi người\dots\ làm anh nhớ tới Honoka-san.}''

Tôi nghiêng đầu.

``\vie{Ý Onii-sama là\dots?}''

``\vie{Ở vòng lặp trước,}'' anh tiếp tục, ``\vie{Honoka-san cũng từng kéo cả nhóm vào một cuộc phiêu lưu chỉ bằng vài câu nói bồng bột như vậy.}''

Onee-sama bật cười khẽ. ``\vie{Đúng thật. Cậu ta bốc đồng đến mức đó cơ đấy. Mặc dù sau đó\dots}''

``\vie{\dots tất cả mọi người, tính cả anh chị, đều bị vướng vào \emph{tai nạn} đó,}'' tôi tiếp lời, không giấu cảm giác lạnh sốn sống lưng khi tôi mơ hồ nhớ lại khoảnh khắc ngôi trường bị đổ sập lúc đó.

Ngay lúc đó, Sakura-san từ bàn gần cửa bước lại. ``\vie{Các cậu đang nói về khu rừng phía sau đền cũ à?}''

Azuki-san gật đầu. ``\vie{Ừ. Bọn tớ đang tính tối nào đó đi thử. Thấy mọi người trong lớp bàn bạc nhau nhiều quá nên cũng tò mò một chút.}''

Sakura-san im lặng vài giây rồi khựng lại như vừa nhớ ra một cái gì đó. Sau một hồi trầm ngâm, cô nói: ``Nếu là khu rừng đó\dots\ tôi cũng đi.''

Tôi không biết tại sao, nhưng giọng của Sakura-san lúc đó rất dứt khoát.

Sau đó Azuki-san bắt đầu lập kế hoạch thật sự.

``Người ta nói nghe thấy tiếng hát lúc khoảng chín giờ tối đúng không?''

``Ừ,'' Honoka-san nói.

``Vậy thì chúng ta cũng đi lúc chín giờ.''

Shion-san vỗ tay tán thành: ``Chủ nhật này nhé!''

``Được đấy,'' Touka-san nói.

Thế là chỉ trong một giờ nghỉ ngắn ngủi, một chuyến đi vào rừng đã được quyết định. Bất ngờ thật.

\begin{center}
  \uline{Kết thúc hồi tưởng.}
\end{center}

Chúng tôi tiếp tục tiếp bước, bước chân vẫn đều đều trên con đường đất nhỏ dẫn vào lòng rừng, như những hành khất đang bước theo tiếng chuông vô hình vang lên từ phía trước.

Tôi lặng lẽ quan sát từng khuôn mặt: người hồ hởi, người căng thẳng, người mải mê cảnh đêm, người thì dán mắt vào bóng tối trước mặt, mỗi biểu cảm như một mảnh ghép riêng trong đoàn hành khất đang bước về phía cánh rừng đen ngòm.

Shion-san và Honoka-san thì cứ tưng tửng trước mặt, thi thoảng châm chọc nhau. ``Tay cậu trông hơi run đấy. Sợ hả?'' cô bạn tóc vàng vỗ vai cô bạn tóc xám, cười tươi.

Shion-san hất cằm: ``Đ-Đây gọi là nhịp điệu thám hiểm, còn ngươi mới là người đang sợ đấy, chú hề.''

``Ồ, thế cậu không run sao?'' Honoka-san nháy mắt. ``Tối nay đừng có bám lấy tớ khóc nhè nhé.''

``Ai bám ai chứ?'' Shion-san đáp lại, giọng đầy thách thức.

Azuki-san đi ở giữa, cười nhẹ: ``Hai cậu thú vị thật.''

Phía trước họ, Saya-san bước chậm, mắt nhìn lên trời đầy xa xăm: ``Nếu đàn tế cổ thật sự từng tồn tại, mảnh đất này có lẽ vẫn còn vang vọng lời nguyện nào đó chưa được cất lên\dots''

Touka-san gật nhẹ, tay vẫn liên tục ghi chép: ``Âm thanh có thể bị cây cối ghi lại, như dữ liệu trong một thí nghiệm dài.''

Saki-san cúi xuống, nhặt cánh hoa dại ven đường, thì thầm trong mơ mộng: ``Ban đêm, nhịp sinh học của chúng đổi khác\dots\ Có thể chúng ta sẽ thấy màu sắc mà chưa ai từng đặt tên.''

Mari-san im lặng, chỉ quan sát tỉ mẩn khung cảnh xung quanh bằng đôi mắt đa sắc. Sakura-san bước song song, mắt hướng vào khoảng tối trước mặt, không nói một lời.

Tôi, Onii-sama và Onee-sama đi sau cùng, lặng lẽ nhìn những người bạn đang đi cùng chúng tôi.

Tôi ghi nhớ mọi chi tiết, cảm nhận nhịp tim mình đập khẽ khi ánh đèn cuối con đường le lói rồi tắt hẳn.

Nhưng trong lòng tôi vẫn còn vất vưởng một câu hỏi: những linh hồn mà tôi đã bắt gặp suốt những ngày qua, chúng muốn gì ở chúng tôi? Và lời nguyền kia, phải chăng vẫn còn đó, chỉ đổi hình thay dạng, len lỏi trong tiếng hát trẻ con giữa rừng đêm?

Có lẽ sau chuyến đi đêm nay, mọi thứ sẽ đựoc sáng tỏ. Dù chỉ là một phần nhỏ trong đó\dots

\ct

Sau một hồi đi bộ, chúng tôi đã tới gần mép rừng. Từ lúc chúng tôi được tái sinh tại đây, vẫn chưa một ai - ngoài Shino-san và Sakura-san - biết được rằng chính chúng tôi là những người đã khôi phục lại cả cánh rừng này.

Con đường đất nhỏ dẫn vào giữa những thân cây tối đen như những cột trụ khổng lồ. Những cơn gió thổi qua tán lá tạo thành âm thanh xào xạc.

Tôi vô thức nhìn sang Sakura-san. Nhưng\dots\ trông cậu ấy có vẻ không bình thường. Cậu ấy đang rất căng thẳng, nhưng không phải kiểu sợ hãi mà chúng tôi thường thấy.

Vai Sakura-san hơi cứng lại. Ánh mắt cô hướng về phía khu rừng như đang nhìn vào một ký ức xa xăm.

Khuôn mặt đó khiến tôi chợt nhớ lại\dots\ vòng lặp quá khứ. Đêm mà nhóm của Honoka-san cùng ba anh em tôi tới trường để tìm tung tích bức tranh. Lúc đó Sakura-san cũng có biểu cảm giống hệt như vậy. Một sự bình tĩnh lạnh lùng, nhưng ẩn dưới là một nỗi sợ rất sâu.

Onee-sama phát hiện ra điều ấy sớm hơn tôi.

``\hl{Sakura-san,}'' chị cất tiếng nhỏ, ``\hl{cậu ổn chứ?}''

Sakura-san khẽ gật. ``Chỉ là\dots\ nơi đây gợi lại vài kỷ niệm.''

Onii-sama hướng mắt về phía rừng, chắp hai tay ra sau: ``Ba linh hồn hôm qua,'' anh chậm rãi kể, ``một ở rạp chiếu phim, một trong con hẻm gần nhà, một tại tiệm lưu niệm.''

Onee-sama khoanh tay trước ngực. ``Chúng chưa gây hại gì cho chúng ta.''

Tôi thì thầm: ``Chúng lại hiện ra bên cạnh những người từng kẹt trong vòng lặp.''

Anh tôi gật nhẹ: ``Game-san cũng gợi ý rằng cội nguồn của chúng có thể nằm trong khu rừng này.''

Sakura-san nhìn sâu vào bóng cây. ``Khi các cậu đã nói rõ thế\dots\ chúng ta càng cần bước vào.''

Trước mặt chúng tôi, những bậc đá rêu phong dẫn lên nằm im lìm dưới ánh trăng nhạt. Phía sau đền, cánh rừng đen kịt trải dài như một bức màn bí ẩn, những tán cây rung rinh trong gió đêm tạo thành âm thanh xào xạc quen thuộc. Cảnh tượng này\dots\ chắc chắn tôi đã từng thấy ở đâu rồi. Không chỉ tôi - tôi có thể cảm nhận được cùng một luồng ký ức đang chảy qua Onii-sama, Onee-sama, và cả Sakura-san. Một cảm giác thân quen đầy mạnh mẽ, như thể chúng tôi đã từng đứng ở chính nơi này, trong một thời điểm nào đó.

\dots Đó chẳng phải là con đường dẫn lên ngôi đền cổ đó sao!? Không thể nhầm lẫn vào đâu được\dots

``\textbf{Tiến lên thôi, mọi người!}''

``\textbf{Đi thôi nào, các chiến hữu của ta!}''

Tiếng kêu gọi của Azuki-san và Shion-san vang lên đưa chúng tôi về thực tại.

Cả nhóm bắt đầu di chuyển về phía lối mòn. Tôi hít một hơi đầy lồng ngực. Ánh đèn đường dần dần bị bóng tối của khu rừng nuốt chửng.

Và, như đã định sẵn mục tiêu trước mặt, chúng tôi bắt đầu bước vào trong rừng.

\ct

Con đường dẫn lên ngôi đền cũ hiện ra trước mắt chúng tôi như một dải tối kéo dài giữa rừng, nhưng trong tôi nó lại như một dải phim cũ đang chạy chậm. Từng bậc đá xanh rêu mờ dưới ánh trăng, chỗ thấp chỗ cao, lúc khúc khuỷu lúc trơn trượt, đều in vào mắt tôi với cảm giác đã từng bước quá nhiều lần đến mức cơ bắp nhớ rõ độ nghiêng. Tôi biết trước chỗ nào sẽ có vết nứt hình con rắn, chỗ nào rêu tròn như vết roi thời gian, chỗ nào bậc thang bị khuyết một mảnh nhỏ đủ để kẹp gót giày. Tất cả đều y hệt, chỉ khác ánh sáng đêm nay pha thêm sương, khiến mọi thứ như bản sao mờ nhạt của một bản gốc từng thuộc lòng.

Cánh rừng hiện ra như một bức tường đen kịt, dày đặc đến mức ánh trăng cũng phải chật vật lọt qua từng kẽ lá. Những thân cây to lớn, cao vút, vỏ xù xì như được đẽo gọt bởi thời gian và gió rét, chen chúc nhau mọc thành hàng, thành lối, như những cột trụ khổng lồ chống đỡ một mái vòm lá rậm rạp. Cành cây đan xen, quấn quýt, tạo thành những mảng tối sẫm không lối thoát. Dây leo buông xuống từng chùm, lúc lắc trong gió như những sợi tóc người chết. Dưới chân, lớp lá khô dày cả gang tay, ẩm ướt, mục nát, phát ra mùi đất pha lẫn mùi xác cây đang thối rữa.

Nhưng tôi biết trước đây, nơi này từng khác hẳn. Như lời mà Saki-san đã từng nói, chỗ này đã từng là một nơi mà hoa màu không thể nào phát triển. Đó là cho đến khi chúng tôi gieo những hạt giống thần kỳ mà Sakura-san đã cất công nghiên cứu, khi đó, nó không chỉ là hoa bỉ ngạn đỏ mà còn cả hy vọng, rừng không còn là rừng—nó là một vườn hoa khổng lồ. Cây cối không chen lấn, mà mọc thành từng hàng thẳng tắp, như được ai đó dịu dàng sắp đặt. Lá cây không dày đặc, mà xòe ra như những chiếc lông vũ, để ánh sáng có thể rọi xuống từng kẽ đất. Hoa nở khắp nơi, không chỉ bỉ ngạn, mà cả những bông hoa trắng tinh, vàng tươi, tím nhạt, mọc thành vòng tròn, thành dải lụa, thành những con đường hoa dẫn vào trái tim rừng. Lúc ấy, rừng không đáng sợ. Lúc ấy, rừng là một bản giao hưởng của sự sống.

Giờ đây, khi màn đêm buông xuống, cánh rừng đầy màu sắc đó bỗng chốc lại im ắng và có gì đó ma mị tới lạ. Lời nguyền của cánh rừng có lẽ vẫn chưa thực sự được hóa giải, nhưng tôi không nghĩ đó là lý do duy nhất.

Tôi liếc sang Onii-sama. Ánh đèn pin của anh cầm hất lên từng bậc đá, nhưng đôi mắt anh không theo dõi ánh sáng mà nhìn xa hơn, như đang đọc một dòng chữ vô hình khắc trên không trung. Onee-sama bước sát bên anh, mũi giày chạm nhẹ vào viên đá thứ mười tính từ dưới lên; chính nơi ấy, ở một vòng lặp trước, chị đã khiêng anh ấy lên cùng với Shino-san. Cả hai không nói gì, nhưng tôi thấy khóe miệng cô mím lại, tay kia nắm chặt vạt áo gió như để ghim mình vào hiện tại, sợ một cơn gió vô hình sẽ cuốn chúng tôi trở lại quá khứ.

Sakura-san đi sau cùng, bước chân cô nhẹ đến mức không nghe tiếng, nhưng tôi cảm nhận rõ mỗi khi cô dừng lại, nơi ánh trăng xuyên qua tán lá tạo thành một hình trăng khuyết trên đá. Cậu ấy ngẩng lên, mắt không nhìn đền mà nhìn vào khoảng không phía trước, nơi ở một thời điểm khác cô từng đứng một mình, tay cầm hạt giống hoa bỉ ngạn, nhưng rồi bệnh tật kéo dài đã khiến cho nhiệm vụ không thành. Tôi biết cô đang đếm lại nhịp tim của chính mình, vì tôi cũng vậy: nhịp tim ấy đập không theo thời gian hiện tại mà theo nhịp của một bản nhạc từng ngân lên từ lòng đất, chỉ những ai từng kẹt trong vòng lặp mới nghe được.

Bốn chúng tôi bước như bống người đang đi trên một bản in lặp đi lặp lại vô tận: mỗi bước chân tái hiện đúng khoảng cách cũ, mỗi cái chạm đá lại gợn lên một lớp ký ức mờ. Tôi cảm thấy mình không chỉ đang leo lên đền mà còn đang leo qua những lớp thời gian chồng chéo, vòng này đè lên vòng kia, tạo thành một xoáy hố đen nhỏ bằng đá rêu, hút tất cả tiếng cười, tiếng khóc, lời nguyền và hơi thở của những đứa trẻ chưa từng được sinh ra. Trên bậc cao nhất, tôi dừng lại, lòng bàn chân đặt đúng chỗ từng in dấu máu của ai đó, máu ấy đã rửa trôi từ lâu, nhưng hơi ấm vẫn còn vương trong đá, chờ một ngày chúng tôi quay lại để tiếp tục bước, tiếp tục lặp, tiếp tục quên, hay cuối cùng sẽ nhớ lại mọi thứ.

Tôi bước chậm lại một chút. Một sự thân quen tới rùng mình đầy bất chợt. Con đường này\dots\ tôi đã từng đi qua rồi. Rất nhiều lần là đằng khác. Khi đó, chỉ có bốn người chúng tôi:

Onii-sama. Onee-sama. Tôi. Và Shino-san.

Shino-san không có ở đây lần này. Nhưng dù vậy, chỉ cần nghĩ tới điều đó thôi cũng khiến lòng tôi chùng xuống trong chốc lát.

Tôi nhìn sang hai bên.

Onii-sama vẫn đi phía trước, dáng vẻ bình tĩnh quen thuộc. Onee-sama bước cạnh anh, thỉnh thoảng quan sát địa hình.

Cả hai chắc chắn cũng đã nhận ra. Không cần một lời nói nào, chúng tôi cũng đủ hiểu.

Chúng tôi chỉ tiếp tục đi về phía trước cùng với đoàn người, quyết định giữ ý nghĩ đó trong đầu.

``\textbf{Mọi người, qua đây này!}'' tiếng gọi bất ngờ vang lên từ phía đầu đoàn.

``Saki-san? Có chuyện gì vậy?'' Azuki-san hướng đèn pin về phía cô bạn tóc vàng tết hai bím, giọng hơi ngạc nhiên. Không cần chờ ai hỏi, cô đã chỉ vào chùm hoa tím li ti và bắt đầu tràng giang đại hải.

``Viola mandshurica!'' Giọng cô như thể vừa gặp lại người thân lâu ngày. ``Sumire đấy, đất ẩm pH chua một chút là chúng sướng ngay.''

``\dots Ờm, là sao?'' Honoka-san đứng trơ ra, gãi đầu gãi tai.

``Người Nhật quen gọi là sumire (スミレ)! Chúng thích nơi đất ẩm vừa phải, pH hơi chua, rồi--'' cậu ấy còn chưa dứt lời, Saki-san đã sang bụi khác, tiếp tục thao thao bất tuyệt: ``Còn đây là hakusan ichige (ハクサンイチゲ)\footnote{Tên khoa học là Anemonastrum narcissiflorum, thuộc họ Mao lương.}! Mọc ở đây thì đúng kiểu bất ngờ đấy!''

Azuki-san cười trừ, nhìn cô bạn đang thao thao bất tuyệt về các loài hoa: ``Saki-san lại lên cơn rồi.''

``Đúng vậy nhỉ,'' Saya-san cười khẽ. ``Cứ mỗi lần nói đến hoa là cậu ấy lại say mê như vậy đó\~{}''

Saki-san vẫn say sưa giảng giải, chẳng để ý ai. Cô lần lượt kể về đất rừng, độ ẩm, pH rồi cách hoa tụ thành từng khóm, miệng không ngơi nghỉ.

Bỗng, giữa câu chuyện, cô khựng lại. ``\dots Hả?''

Ánh mắt cô dán vào khoảnh đất trống phía trước, nơi những bông hoa trắng muốt đang rung rinh dưới ánh trăng lạt.

``Bỉ ngạn trắng\dots'' tiếng thì thào thoát ra từ cổ họng cô, nhỏ đến mức nghe như tiếng lá rơi.

Ba chữ ấy đủ khiến tôi, Onii-sama và Onee-sama đồng loạt cứng đờ. Không phải màu đỏ kinh hoàng từng tràn ngập hành lang lớp học, nhưng vẫn là bỉ ngạn; loài hoa nở trên ranh giới, báo hiệu điều gì đó vừa chết, vừa chuẩn bị sống lại.

Vậy mà sao thứ này lại\dots

Saki-san lẩm nhẩm: ``Người ta hay nhớ bỉ ngạn đỏ, nhưng bạch bỉ ngạn cũng chẳng vô hại. Chúng thích bám trên nền mộ cổ, hút chất vôi từ xương\dots'' Cô vẽ một vòng tay quanh khoảnh đất, mắt không rời những cánh hoa mong manh. ``Nhìn cách chúng mọc thành vòng tròn khép kín này đi, không giống tự nhiên chút nào.''

Không khí nặng trĩu. Mari-san khẽ hỏi: ``Cậu muốn nói\dots\ từng có nghĩa trang ở đây?''

Saki-san lắc đầu, nhưng giọng cô khàn đặc: ``Mình không dám chắc. Nhưng nếu có, thì chúng ta đang đứng trên\dots\ một khuôn mộ khổng lồ.''

Tôi liếc về phía Sakura-san. Cậu ấy đứng thẳng băng, hai bàn tay siết chặt đến mức các khớp trắng bệch; ánh mắt cậu ấy như đang nhìn xuyên qua màn đêm, xuyên qua những kẽ lá, để chạm vào thứ gì đó chỉ cô mới thấy. Không một tiếng thở dài, không một lời than, mà chỉ có im lặng nặng trĩu đè xuống vai cậu ấy.

Gió thốc qua khe cổ áo tôi, mang theo mùi đất ẩm pha lẫn hương hoa bỉ ngạn ngọt ngào nhưng lạnh lẽo, như bàn tay ai đó vừa vuốt ve vừa siết cổ. Tiếng lá xào xạc bỗng trở nên rõ ràng hơn bao giờ hết, gần như đang thì thầm bên tai, đe dọa đừng bước thêm nữa. Mọi người đứng thành vòng, ánh đèn pin loang loáng run rẩy trên những bông hoa trắng, khiến chúng trông như những chiếc mặt nạ nhỏ nhắn đang nở nụ cười gượng gạo. Tôi nuốt khan, cảm thấy không khí dày đặc đến mức phải dùng vai để đẩy ra mới có thể bước tới.

``Khoan đã.''

Giọng Touka-san cất lên, đủ lớn để cắt đứt cả mớ âm thanh rợn người đang bám lấy chúng tôi. Cô đưa đèn pin quét chậm từ gốc cây này sang gốc cây khác, như thể đang đo khoảng cách giữa các vì sao. ``Cậu vừa nói `vòng quanh', đúng không?''

Saki-san gật nhẹ, hai bím tóc rung rung trước trán.

Touka-san dịch chân thêm một bước, đặt bàn chân lên mép ngoài vòng hoa trắng như thể đang đặt dấu chấm hỏi lên mặt đất. Đôi mắt cô nheo lại, đầu nghiêng nhẹ, toát lên vẻ của người vừa bắt được một manh mối nhỏ xíu nhưng đủ làm thay đổi toàn bộ bức tranh đêm nay. ``Nhìn kỹ đi,'' cô khẽ nói, giọng trầm xuống như tiếng gió thổi qua kẽ lá, ``rừng ở đây có kết cấu hơi cong.''

Aku-san chớp mắt, còn chưa kịp hiểu chuyện gì đang xảy ra: ``Hả?''

Không đợi ai lên tiếng, Touka-san cúi xuống nhặt một cành cây khô, ngồi thụp xuống đất và bắt đầu vạch những đường ngoằn ngoèo trên nền lá ẩm. ``Nhìn này,'' cô nói, tay liên tục di chuyển, vẽ một cung tròn lớn ôm lấy khoảnh đất trống. ``Từ lúc rẽ vào đường mòn, mình đã thấy lối đi không thẳng mà hơi lượn.'' Cô dừng lại, chỉ vào hai điểm đầu và cuối cung. ``Nó không đi thẳng, mà ôm vòng như lưỡi liềm.''

Saya-san nghiêng đầu, mắt dán vào hình vẽ: ``Ý cậu là\dots\ cả khu rừng này có hình vòng cung?''

Cô nàng tóc xanh lam gật nhẹ, mái tóc lòa xòa trước trán. ``Mình từng đọc trong một cuốn sách về kiến trúc nghi lễ cổ,'' cô ngẩng lên, ánh mắt sáng lên như vừa tìm ra mảnh ghép còn thiếu. ``Nếu một khu rừng được giữ nguyên dạng vòng cung như thế này, thì trung tâm của cung thường là nơi diễn ra nghi lễ.''

Onii-sama cất tiếng, giọng trầm ấm vang nhẹ trong đêm. ``Ở đâu ra cuốn sách đó?''

``Thư viện thành phố,'' Touka-san đáp ngay, không chút do dự. ``Mình nhớ nó có đề cập đến rừng vòng cung được dùng làm đài tế thời Nara từ khoảng 1300 năm về trước,'' cô nàng tiếp tục vẽ, cành cây khô kêu lạo xạo trên mặt đất. ``Nếu đúng như vậy, chỗ chúng ta đang đứng chỉ là vành ngoài. Trung tâm nghi lễ phải ở sâu hơn, nơi nào đó giữa lòng rừng.''

Một khoảng lặng ngắn bao trùm. Gió thốc qua khe lá, mang theo mùi đất ẩm và hương hoa bỉ ngạn thoang thoảng. Onii-sama cúi xuống nhìn bản phác thảo, khóe môi khẽ nhếch như thể đang tính toán điều gì.

Onee-sama khoanh tay, mắt dán vào điểm giữa cung tròn. ``Cậu nghĩ trung tâm ở đâu?'' chị hỏi, giọng khẽ như sợ đánh thức thứ gì đang ngủ dưới lớp đất mùn.

Touka-san nhấc cành cây lên, gõ nhẹ vào điểm chính giữa hình vẽ. ``Khoảng\dots\ đâu đó quanh đây,'' cậu ấy thì thầm, như thể chính bản thân cũng chưa dám tin vào lời mình vừa nói. Gió rừng lúc ấy chợt lạnh hơn, thổi tung tóc cô, cũng thổi cả mùi đất ẩm vào mũi mọi người.

Một sự im lặng bám lấy chúng tôi thêm vài nhịp tim. Rồi, từ phía sau vòng người, một giọng nói quen thuộc vang lên, chậm rãi như tiếng chuông gió đêm.

``Nếu vậy thì\dots\ thứ này có thể giúp chúng ta.''

Saya-san bước lên. Cô nàng không vội, chỉ lần mở dây túi vải nhỏ treo bên hông. Từ trong đó, cô rút ra một tờ giấy đã ngả màu, gấp bốn vuông vức. Những nếp gấp cũ kỹ in hằn vết thời gian, nhưng vẫn còn mùi mực khô thoang thoảng.

Tôi nghiêng đầu, chưa kịp hỏi, cô nàng đã tự trả lời: ``Bản đồ đàn tế cổ đó\~{} Mình tìm được trong thư viện trường.''

Cậu ấy trải tờ giấy ra trên nền lá. Ánh đèn pin hắt lên, lộ ra những đường mực nhạt màu chì, vẽ khu đất hình bán nguyệt. Rải rác trên đó là những ký hiệu tròn, tam giác, lưỡi liềm tưởng chừng vô nghĩa, nhưng lại sắp xếp theo trật tự rõ ràng. Tôi lập tức nhận ra ngay: nó trùng khít với hình vẽ Touka-san vừa hoàn thành.

Aku-san cau mày: ``Thư viện nào lại cất bản đồ kiểu này?''

Saya-san chỉ mỉm cười, không đáp. Cô nàng chỉ tay lên tờ giấy, bắt đầu thuyết minh từng chấm sao: ``Đây là điểm hướng trăng tròn vào đêm thu phân. Đây là trục Bắc Cực cổ. Còn đây,'' cô dừng lại, đặt ngón tay vào trung tâm bản đồ, ``là nơi cử hành nghi lễ chính.''

Touka-san ngẩng lên, mắt sáng lên như vừa tìm được mảnh ghép cuối. Cô đặt cành cây xuống, so hai bản vẽ, rồi thốt lên nhẹ: ``Kh-Khớp y hệt kìa!''

Saya-san gật đầu, giọng vẫn đều đều: ``Ngôi đền chúng ta thấy chỉ là điểm phụ. Trái tim thật sự của cấu trúc\dots\ sẽ nằm đâu đó ở sâu hơn trong rừng.''

Gió đêm lướt qua, xoáy lá rừng lên thành từng đợt sóng thì thầm. Mọi người đứng im như thể vừa bước vào một bức tranh thủy mặc chưa khô mực; ánh đèn pin run rẩy, chỉ dám dừng trên mép vòng hoa bỉ ngạn như sợ đánh thức điều gì đang nằm dưới đó. Tôi nghe rõ tiếng tim của mọi người đập thình thịch trong thinh không nặng trĩu, không phải chỉ riêng mình tôi.

Sakura-san khẽ lùi nửa bước, vai chạm phải vai tôi. ``\hl{Cậu cũng thấy đấy chứ?}'' giọng cậu ấy chỉ đủ hai chúng tôi nghe: ``\hl{Bản đồ, hoa, lối đi\dots{} ba mảnh ghép này không thể nào khớp hơn được nữa.}''

Tôi gật, không dám nói to vì sợ âm thanh sẽ làm vỡ cái màng mỏng manh đang bao bọc chúng tôi. ``\hl{Như kiểu\dots{} ai đó đã dựng sẵn một cái bẫy cho chúng ta,}'' tôi thì thào, ``\hl{và chúng ta đang đi theo hướng của nó.}''

Onee-sama ôm cánh tay mình, mắt vẫn dán vào tấm bản đồ cũ. ``Nếu đây là vòng lặp cũ,'' chị tôi nói không thành lời, ``thì lần trước\dots\ chúng ta đã từng tới tận trung tâm sao?'' Không ai trả lời, nhưng tôi biết cả bốn người chúng tôi đang cùng nhớ ra một đêm mưa rơi lất phất, tiếng chuông gió đền cổ và mùi hoa bỉ ngạn như bây giờ. Ký ức không lời lẫn vào hơi sương, bám trên làn da lạnh ngắt.

Mari-san cười gượng, phá vỡ im lặng: ``Chắc mình từng đi học ngoại khóa ở đây\dots\ hay là lúc đi dạo vẽ phong cảnh?'' giọng cô rung nhẹ, không đủ thuyết phục chính mình. Tôi liếc sang Onii-sama; anh khẽ lắc đầu, ý bảo ``đừng nói gì thêm.'' Onee-sama siết chặt tay áo tôi, như thể sợ tôi sẽ bước lệch khỏi vòng tròn ánh sáng mong manh.

Sakura-san đứng im như tượng đá, đôi mắt lướt nhanh từ vòng hoa bỉ ngạn trắng muốt, sang bản phác thảo đất rừng của Touka-san, rồi dừng lại ở tấm bản đồ cổ mờ mực trong tay Saya-san. Ba mảnh ghép ấy chồng lên nhau, hiện ra trong đầu cô thành một bức tranh quen thuộc đến rợn người. Cô khẽ mấp máy môi, giọng trầm ấm nhưng run nhẹ: ``Mạch ký ức đã thông, mọi thứ đang trở về đúng vị trí cũ.''

Cô quay sang ba anh em tôi, ánh mắt sâu thẳm: ``Đây không phải cuộc dạo chơi đêm khuya. Đây là dấu vết một vòng lặp vẫn hở mạch, chờ chúng ta bước vào để khép lại.'' Cô ngẩng lên nhìn tán cây đen kịt, tiếp tục thì thầm như sợ đánh thức thứ gì đang ngủ: ``Chúng ta đang giữa tấm mạng nhện thời gian; mỗi bước chân rung một sợi chỉ, và ở tâm mạng có kẻ đang nắm chuông báo thức.''

Gió chợt rít dài, cuốn lá khô thành vòng xoáy lạnh lẽo quanh bờ giày. Đêm nay không chỉ là cuộc truy tìm tiếng hát trẻ con, mà là bước ngược vào ký ức tập thể, thứ chúng tôi đã vô tình thả lạc ở một vòng lặp trước. Phía trước, con đường đất nhỏ như dải băng mờ ảo dẫn thẳng vào bụng tối, nơi ánh đèn pin cũng chỉ còn là hạt sương le lói.

Azuki-san quay lại, nụ cười vừa hồi hộp vừa nôn nao: ``Tiếp tục thôi chứ?'' Im lặng đồng ý. Không ai thốt lời, nhưng tất cả đồng loạt bước - sâu hơn, sâu nữa - để hòa mình vào khu rừng đang mở toàn bộ lá phổi đen đặc hít lấy chúng tôi.

\ct

Chúng tôi lại lầm lũi bước sâu hơn vào màn đêm rừng rậm. Gió cuối xuân lùa qua kẽ lá, mang theo hơi lạnh bám vào da thịt, khiến tôi không khỏi rùng mình. Ánh đèn pin loang loáng trên những thân cây to lớn, in bóng chúng thành những cột trụ đen ngòm vô tận.

Có bản đồ cổ của Saya-san cùng bản phác thảo cung tròn của Touka-san (được vẽ lại trên giấy), cả nhóm giờ đã xác định được phương hướng rõ ràng. Chúng tôi bắt đầu tiến chậm rãi, dò từng bước như người mù đang đi trên dây. Mỗi lần dừng lại, cô bạn Thần đồng lại cúi xuống, dùng ngón tay ấn nhẹ lớp đất, đo độ nghiêng rồi gật gù. Cô bạn Chiêm tinh thì so từng tàng cây, tảng đá với nét mực nâu nhạt trên tờ giấy, môi mấp máy đếm bước. Còn Saki-san\dots\ cô nàng vẫn không bỏ được thói quen ``bổ sung'' thực đơn thực vật cho bộ sưu tập, thỉnh thoảng reo lên khi phát hiện một loài hoa lạ mọc lẫn trong bụi rậm.

Đường rừng không dài nhưng quanh co, đôi lúc tưởng đã lạc, may nhờ có hai bản đồ soi đường, cả nhóm lại tìm ra hướng đi. Lá khô dưới chân giòn tan, tiếng vỡ lách tách vang lên giữa bầu không khí tĩnh lặng đến nghẹt thở. Mỗi người đều cảm nhận rõ mùi đất pha lẫn hương hoa bạch bỉ ngạn còn vương trong gió, như lời nhắc nhở chúng tôi đang bước vào một vùng đất không dành cho kẻ sống.

Rồi, sau cái khúc cua gấp cuối cùng, mái đền cổ hiện ra. Nó nằm lặng lẽ giữa khoảng trống tròn trịa, tựa hồ đang chờ đợi chúng tôi từ ngàn xưa. Ánh trăng lúc này đã lên cao, phủ lên ngói rêu một lớp bạc mờ ảo. Tôi rút điện thoại, màn hình hiển thị 21:48. Không ai nói gì, nhưng tôi biết, tất cả đều thầm đếm từng nhịp tim trong bóng tối đang dày đặc bao quanh.

Touka-san khép tấm bản đồ cổ lại, ánh mắt cô dừng trước mái đền lợp rêu phong ẩn hiện trong sương khuya. ``Chính tại đây,'' cô thì thầm như sợ đánh thức ai, ``nghi lễ đã từng được khai mở.''

\dots Khoan đã.

Tôi bất giác thắt tim. Ngôi đền kia\dots\ không lẽ lại là nơi ấy?

Những chiếc đèn lồng bằng giấy bồi treo dưới hiên gỗ tỏa ánh vàng vọt, lúc tỏ lúc mờ theo nhịp gió thổi từ rừng sâu. Ánh sáng mong manh như hơi thở cuối cùng của ký ức, vừa chực tắt ngấm vào màn sương dày.

Toàn bộ ngôi đền im lìm, đứng đó từ ngàn đời, dường như đã quá quen với việc làm chứng cho những điều không bao giờ được kể hết.

Tôi không thể rời mắt. Trong một vòng lặp trước - khi cánh rừng này từng tràn ngập sức hoa từ hạt giống do Sakura-san tạo ra - chúng tôi cũng đã đứng trên bậc đá này, tin rằng mọi thứ đã kết thúc tốt đẹp. Giờ đây, cái kết ấy như bị gió mạnh giật lại, chỉ còn lại hơi lạnh quanh cổ chân và mùi rêu ẩm ướt bốc lên từ kẽ đá.

Tôi thở dài, nhẹ đến mức chính mình cũng không nghe thấy. Có lẽ, vòng lặp chưa bao giờ thật sự đóng lại; nó chỉ khép mình trong bóng tối, chờ một ngày chúng tôi quay lại, như người thợ giăng mắc cỡ mới.

``Ồ\dots''

Tiếng nhẹ nhàng thoát ra từ Azuki-san như một làn khói trắng lẫn trong sương đêm. Ánh mắt cậu ấy lướt nhanh qua mái đền cổ, rồi dừng lại trên những bậc đá rêu phong, như thể vừa tìm thấy một bức tranh thời thơ ấu bị bỏ quên giữa rừng sâu. ``Không ngờ trong lòng rừng lại giấu một ngôi đền như thế này,'' giọng của cô nàng ngập tràn niềm khám phá.

Shion-san khoanh tay trước ngực, nheo mắt. Gió làm mái tóc bạc xám bay nhẹ, che khuất một phần khuôn mặt, khiến nụ cười thoáng hiện trở nên bí ẩn. Giọng cạua ấy trầm xuống, như tiếng trống thúc nhịp chậm trong đêm không trăng:``Hừm. Đúng là nơi thích hợp cho những nghi thức cổ đại.''

Aku-san lập tức nhăn mặt, vẻ mệt mỏi quen thuộc. ``Cậu lại bắt đầu thể hiện rồi đó, Shion-chan\dots'' Nhưng cô bạn thân chỉ mỉm cười, nửa miệng, nửa mắt. ``Biết đâu đấy, Aku. Có khi nơi này từng là nơi phong ấn linh hồn chẳng hạn.''

Câu nói vừa dứt, Aku-san đã rít lên, giọng như bị ai véo vào tai: ``Đừng có nói mấy thứ đáng sợ như vậy!'' Azuki-san bật cười, tiếng cười lanh lảnh làm vỡ đi bầu không khí căng thẳng chỉ vừa kịp hình thành.

Phía sau ba người họ, Saya-san, Touka-san và Saki-san đứng im như những bức tượng nhỏ. Ánh mắt họ dán vào ngôi đền, nhưng dường như họ không thấy ngạc nhiên, có lẽ vì họ đã từng đứng ở đây từ trước đó, nếu không vì mục đích tâm linh cầu Thần thì cũng là mục đích nghiên cứu khoa học nào đó khác.

Honoka-san và Mari-san thì khác. Họ lần đầu đặt chân tới nơi này. Honoka-san mở to mắt, môi hé mở, thốt ra tiếng ``Wow\dots'' nhỏ xíu, như sợ lỡ mất một giai điệu bí ẩn nào đó đang lẩn khuất giữa những kẽ lá. Mari-san nhìn quanh, ánh mắt tò mò, như thể đang cố ghi nhớ từng chi tiết để vẽ lại trên giấy khi trở về.

Tôi bước chậm lại, gần hơn bậc đá. Mỗi bước như đạp lên một mảnh ký ức. Cảm giác dưới chân lạnh, cứng, hơi nhám\dots\ giống hệt lần trước. Không cần nhìn, tôi cũng biết chỗ nào có vết nứt hình con rắn, chỗ nào có mảng rêu tròn như vết roi khắc thời gian. Có những nơi, dù chỉ đến một lần, lại in sâu vào máu thịt như một bản nhạc buồn không bao giờ phai. Và nơi này chính là một trong những bản nhạc ấy.

Onee-sama chợt quay nửa người, ánh mắt sắc như lưỡi dao mổ đêm. ``Azuki-san?''

Cô bạn tóc đen pha hồng ngoảnh đầu lại. ``Sao thế?''

``Cậu kể lại tin đồn trong trường lúc trước đi.''

``Tin đồn\dots\ à, cái mà mọi người xì xào ở mấy anh chị lớp trên đấy hả?'' Honoka-san chêm vào.

Onee-sama gật, không nói gì thêm. Không khí quanh bậc đá như cứng lại, chỉ còn tiếng gió rít của buổi đêm cuối xuân lành lạnh.

Azuki-san đưa ngón tay lên cằm, suy tư một hồi rồi bắt đầu kể: ``\hl{Nghe đâu, nếu can đảm bước vào rừng đúng lúc giờ Hợi chính khắc, ta sẽ thấy lũ trẻ ngồi thành vòng\dots{} quay lưng về phía mình. Chúng không bao giờ quay đầu lại, dù gọi thế nào.}''

``Giờ Hợi chính khắc?'' Onii-sama nghiêng đầu.

``Tớ cũng không hiểu nữa\dots'' cô bạn bờm cáo gãi đầu gãi tai. ``Chỉ biết các anh chị lớp trên bảo thế.''

Onii-sama khẽ nhếch môi, ánh mắt như đang lật từng trang sách cũ trong đầu. ``Giờ Hợi chính khắc\dots\ Nếu xét theo lịch Can Chi, Hợi ở vị trí cuối cùng theo Can Chi.'' Anh vẽ một vòng tròn nhỏ trên không trung bằng ngón tay, như thể đang vẽ một đồ thị hàm số. ``Nhưng điều quan trọng là cách tính giờ.''

Touka-san gật nhẹ, mắt sáng lên như vừa tìm được mảnh ghép quan trọng của một bài toán khó. ``Đúng vậy. Trong hệ thống Can Chi, một ngày được chia thành 12 canh, mỗi canh tương ứng với 2 giờ hiện đại. Thìn là canh thứ 5, tính từ Tý (23h-1h).'' Cậu ấy đưa tay đếm từng ngón, giọng như máy tính đang chạy: ``Tý (23-1), Sửu (1-3), Dần (3-5), Mão (5-7), Thìn (7-9)\dots\ và cứ như vậy.''

``Nhưng Honoka-san nói là `giờ Hợi chính khắc' mà?'' Aku-san chau mày, vẻ mặt như đang cố theo kịp một bài giảng toán học.

``Cứ để hai người họ giải thích xem nào,'' Azuki-san lên tiếng, đèn rọi vào bản vẽ ở dưới mặt đất.

Onii-sama mỉm cười, tay chỉ vào thứ Touka-san vừa vẽ. ``Vì đây là \textbf{giờ Hợi chính khắc}, không phải giờ Hợi thuần túy. Khắc thì khác một chút, mỗi giờ con giáp đó lại chia thành 8 khắc, mỗi khắc là 15 phút.'' Anh khẽ gõ trán, như thể đang trích xuất dữ liệu từ một thuật toán phức tạp. ``Chính khắc nằm ở chính giữa, tức là khắc thứ tư, tức một giờ.''

``Vậy nghĩa là 9 giờ cộng thêm 1 giờ nữa sẽ là 10 giờ đêm,'' Sakura-san kết luận, giọng như vừa giải xong một bài toán hóc búa. ``Giờ Hợi chính khắc là 22:00 hiện đại - thời điểm mà lũ trẻ trong tin đồn sẽ xuất hiện, theo những gì mà Honoka nói.''

Một khoảng lặng ngắn bao trùm. Gió thốc qua khe lá, mang theo mùi đất ẩm và hương hoa bỉ ngạn thoang thoảng.

``Sắp đến giờ rồi,'' Azuki-san nói, giọng khẽ như sợ đánh thức ai, ``giờ Hợi chính khắc\dots''

Shion-san mỉm cười, mắt nheo lại: ``Cửa đã mở. Ta muốn xem đối thủ ở đây mặt mũi như thế nào!''

Aku-san nuốt khan, nắm chặt tay áo Shion: ``Đ-Đừng có nói thế chứ!''

Gió thổi qua khe cổ áo, lạnh như hơi thở từ miệng hầm cũ. Tôi thấy anh mình khẽ giơ tay, mặt đồng hồ lấp lánh dưới ánh đèn pin: đúng 22:00, kim giây dừng lại một nhịp, như có ai vừa gõ cửa thời gian.

*VÙÙÙ*

Gió chợt thổi mạnh, không còn là làn hơi mơn man mà thành trận lốc nhỏ xé rừng. Lá đồng loạt run rẩy, xào xạc thành âm thanh có nhịp, như có ai đang gõ nhịp bản trường ca trên mặt trống cây.

Tôi nín thở, nhận ra nhịp điệu ấy không hề tự nhiên; nó tròn vành, trầm bổng, lúc lên lúc xuống, đúng kiểu giai điệu ru con mà người mẹ thầm thì trong đêm không trăng. Tiếng hát của gió, hay gió đang bê đôi môi ai đó hát, len lỏi vào tai, mềm như lụa nhưng lạnh như băng.

Shion-san, vừa còn cười cợt chuyện phong ấn linh hồn, bỗng bặt tăm như bị bịt miệng. Cả nhóm đồng loạt khựng lại, tiếng bước chân tắt ngấm, tiếng lá lúc này thành thứ duy nhất còn vang. Không khí hạ xuân bỗng như trượt thẳng vào đông; tôi rùng mình, cảm giác có dải lụa ướt quét dọc sống lưng, kéo từ gáy xuống tận cổ chân. Ánh đèn pin run rẩy, không vì gió mà vì bàn tay người cầm đang siết chặt đến trắng đốt ngón.

Và sau đó\dots\ chúng đã thật sự đến.

Chúng đến thật rồi. Không nhảy xổ, không ập đến như bóng ma trong phim, mà từ tốn như thủy triều lên. Từng chấm đen nhỏ xíu mở ra nơi mép đất trống, lúc đầu chỉ là vài mảnh mờ, sau thưa thớt dày dần, như ai rắc hạt mè đen xuống tấm vải mờ ảo. Những bóng nhỏ ấy không đi, không chạy, chỉ lặng lẽ hiện ra, đứng thành vòng cung ôm lấy chúng tôi, như đợi một lệnh vô hình để cất bước.

Một bóng. Hai bóng. Ba bóng. Rồi nhiều hơn, từng bóng, từng bóng một.

Chúng ngồi thụp xuống đất, đầu gối ép sát ngực, như những búp bê bị bỏ quên. Không ai nhúc nhích, không ai thở mạnh; chỉ có làn gió lạnh luồn qua, làm rung nhẹ chiếc đèn lồng giấy cũ kỹ treo trước hiên đền. Trong vòng tròn âm u ấy, những bóng trẻ chỉ còn là mảng mờ mỏng, mặt mũi nhòa như bức tranh thủy ẩm bị nước mưa dội lâu ngày. Tôi bỗng có cảm giác chúng đang lắng nghe một lời gọi vô hình, chờ một mệnh lệnh từ thế giới bên kia để biết có nên hiện hình hay tiếp tục ẩn mình trong tối.

Giữa đám sương đen đặc quánh đó, một bé gái bước ra-- không, phải nói là ``lộ ra'', vì cô bé không hề cử động chân. Em chỉ đơn giản là hiện diện, như thể màn đêm vừa nhấc một lớp voan mỏng lên để lộ khuôn mặt thật. Khoảng tám tuổi, tóc đen tết đuôi sam thấp thoáng dưới vành nón cổ, vạt kimono nâu sẫn thời Nara phủ kín mắt cá, viền áo thêu họa tiết hoa cúc mất màu.

Điều khiến tôi cứng họng không phải y phục, cũng chẳng phải vẻ cổ xưa toát ra từ em, mà một cái gì đó không giống như những linh hồn khác: đôi mắt xanh lục như hồ nước ngọt giữa rừng sâu, trong veo, bình thản, và lặng lẽ gắn thẳng vào mắt tôi, như thể em vừa nhận ra kẻ lạc bước lọt vào cuộc hẹn ngàn năm trước. Trong khoảnh khắc ấy, tôi hiểu: em không chỉ đứng ngoài vòng tròn, em đứng ngoài cả thời gian, và ánh nhìn ấy chính là lời thông báo rằng màn đêm đã sẵn sàng mở ra điều gì đó không thể quay đầu.

Cô bé đứng yên, lặng lẽ như bức tượng phủ sương, chỉ có hơi thở mỏng manh vừa đủ khiến lá gần đó rung nhẹ; dường như em đang thử thách chúng tôi ai sẽ là người động đậy trước.

Onii-sama bỗng nắm chặt cán đèn pin, khẽ lùi nửa bước: “Hình như\dots\ anh đã gặp cô bé này rồi.” Onee-sama đứng bên cạnh cũng sững sờ, đôi mắt nheo lại: ``Tôi cũng vậy, ngày hôm qua ở rạp chiếu phim, tôi thấy em bé này ngồi một mình hàng ghế cuối.''

Tôi cảm thấy máu chạy rần rần trong thái dương: ``Và em\dots\ trong tiệm lưu niệm bên phố, em đứng trước kệ búp bê. Lúc đó, em ấy bất chợt xuất hiện, nhưng\dots\ sao có thể\dots''

Ba chúng tôi đồng loạt nhìn nhau, tiếng gió như ngừng thở: ``Vậy\dots\ hóa ra chúng ta đều nhìn thấy cùng một người sao?''

Trong lúc chúng tôi còn bị động khi gặp lại cô bé ấy theo cách khó ngờ nhất, những người còn lại cũng dần nhận ra không khí đang cứng lại. Azuki-san giật mình, môi run rẩy: ``Thật\dots\ thật sự có\dots'' Câu nói dở dang tan trong cổ họng, cô nàng lùi nửa bước, bàn tay vô thức tìm kiếm chỗ dựa.

Shion-san, dù mặt đã trắng bệch, vẫn khoanh tay trước ngực như muốn giữ vẻ bình tĩnh thường ngày; tiếng ``\dots ma thật kìa\dots'' thoát ra như một cái gật đầu bất đắc dĩ với thế giới siêu nhiên.

Aku-san không nói gì, chỉ siết chặt tay cô bạn thân, đầu gối khẽ run, mắt không rời bóng hình nhỏ bé kia, sợ rằng nếu chớp mắt, em sẽ biến thành thứ gì đó khủng khiếp hơn.

Touka-san đứng sững, logic trong đầu như bị gió đêm thổi bay. ``Không thể nào\dots'' cậu ấy lẩm nhẩm, ``mọi phân tích của mình\dots\ không lẽ lại\dots?'' Cành cây khô trong tay rơi xuống, tiếng động nhỏ nhưng đủ đánh thức cả nhóm khỏi mê hồn.

Saki-san cúi xuống, mắt dán vào vòng hoa bỉ ngạn trắng như thể chỉ cần nhìn kỹ hơn, cô sẽ tìm ra lời giải thích khoa học cho hiện tượng không-tưởng này; nhưng cánh hoa chỉ rung nhẹ, im lặng giữ bí mật.

Mari-san khẽ mấp máy, tiếng thì thầm mỏng như sợi tơ, tôi chỉ kịp bắt được vài chữ ``\dots có thật\dots'', rồi cô ôm vai mình, lùi vào bóng tối của tán cây.

Saya-san ngẩng cao đầu, ánh mắt lướt trên bầu trời đêm như đang đối chiếu tọa độ các vì sao với bản đồ cổ trong túi; có lẽ cô tìm kiếm một chòm sao nào đó để xác nhận rằng chúng tôi vẫn còn ở đúng thế giới này.

Trong khi cả nhóm còn đứng chôn chân, bốn chúng tôi đã bước lên phía trước, như thể có một sợi dây vô hình kéo họ đến gần hồn ma cô bé mắt xanh lục. Ánh đèn pin run rẩy, chiếu vào khuôn mặt nhỏ nhắn của em, hiện rõ một gương mặt không một chút sợ hãi, chỉ toát lên vẻ buồn xa xăm.

Onii-sama không vội tiến gần. Anh dừng lại, đôi mắt nheo lại, quét từng bóng trẻ như đang đọc một bản đồ ký ức. Khoảng cách giữa họ được anh đo bằng bước chân, bằng nhịp thở, bằng cả sự im lặng đang căng ra giữa người và ma. Rồi, anh khẽ lên tiếng, giọng trầm đến mức gần như tan trong gió:

``Đúng như Game-san dự đoán\dots\ chúng đang bị mắc kẹt ở đây.''

Onee-sama không nói gì. Chỉ một bước chân nhẹ, chị đã đứng chệch trước Sakura-san và tôi; không hoàn toàn che chắn, nhưng đủ để ánh mắt cô nói thay lời: ``Đừng động đến em ấy.'' Đôi mắt chị tôi lướt qua từng linh hồn, như thể đang đếm những mảnh vỡ của một thời đại đã mất.

``Chúng không còn bám theo từng người nữa,'' cô thì thầm. ``Đây thực sự là nơi\dots\ mà mọi thứ tụ tập lại.''

Tôi bước lên. Không nhanh, không chậm, chỉ đủ để không làm vỡ màn đêm đang treo lơ lửng trên khoảng đất trống. Tôi quỳ xuống, đầu gối chạm lá ẩm, đôi mắt ngang tầm với cô bé đứng đầu.

Giọng tôi dịu đến mức gần như không thành lời: ``Chị không đến để đuổi các em đi.'' Cặp mắt hai màu của tôi nói thay lời muốn nói rằng: chị sẽ không làm hại tới các em đâu.

Không có một sự phản hồi, cũng không cử động. Chỉ có đôi mắt xanh lục của cô bé - trong veo, như hồ nước ngừng thời gian - chuyển hướng, nhìn thẳng vào tôi. Như thể có một sợi chỉ vô hình giật nhẹ, em nghiêng đầu. Một cử chỉ nhỏ, nhưng đủ để tôi biết: em đang lắng nghe đây.

Sakura-san bước lên. Cậu ấy không nhìn tôi, cũng không nhìn cô bé, mà nhìn quanh, như đang đọc một bản thánh ca cũ bằng mắt. Hoa bỉ ngạn, ngôi đền, vòng tròn đất trống\dots\ tất cả như lùi lại một bước để cô nhìn rõ hơn.

Cậu ấy khẽi, giọng như tiếng lá rơi: ``Đây là nơi nghi thức thật sự đã bắt đầu.''

Câu nói vừa dứt, không khí rung lên. Một sự dao động mạnh rõ ràng. Những linh hồn như làn sương bị khuấy động, mờ dần rồi lại hiện rõ, như thể có ai vừa chỉnh lại tiêu cự của một bức ảnh cũ. Cô bé mắt xanh không nhìn tôi nữa. Em nhìn Sakur-san lâu hơn, sâu hơn, như thể đang nhận ra một khuôn mặt từng bị lãng quên trong một giấc mơ dài.

Tôi cảm thấy nó, một sự công nhận. Không bằng lời, không bằng tiếng, chỉ bằng cách không khí giữa họ trở nên mỏng manh đến mức có thể cắt bằng tay.

Một sự công nhận rằng những linh hồn này đã từng là những cô bé, cậu bé.

Một sự công nhận cho thấy sẽ có những người sẵn sáng lắng nghe những điều những linh hồn này chia sẻ.

Và, một sự công nhận rằng sẽ có người tói để cứu rỗi những linh hồn này về nơi an nghỉ cuối cùng.

Chỉ là\dots\ có vẻ như các em ấy vẫn chưa sẵn sàng để mở lời. Chúng chỉ hiện về, lặng lẽ khẳng định rằng: ``Chúng em vẫn còn ở đây.''

Im lặng.

Onii-sama lên tiếng, không lớn, nhưng đủ để cả khu rừng phải lắng nghe:

``Nếu các em muốn được lắng nghe\dots\ bọn anh sẽ quay lại.''

Một cơn gió mạnh bất ngờ thổi qua. Không phải từ rừng — mà từ *bên trong* khoảng đất trống. Đèn lồng rung lên, tiếng gỗ kêu ken két như tiếng cửa sổ cũ vừa mở ra sau ngàn năm. Những linh hồn, từng bóng một, bắt đầu tan, không phải biến mất, mà trở về, tựa như chúng vừa được tha thứ một lỗi họ không bao giờ kịp nói ra.

Chỉ trong vài giây, khoảng đất trước ngôi đền trở nên trống rỗng. Không còn bóng trẻ, không còn ánh mắt xanh lục, chỉ còn tiếng gió rít qua khe lá, như một bản nhạc nền của một vở kịch đã khép màn.

Nhưng tôi biết, chúng chưa đi xa. Chúng chỉ lùi lại, và chờ một thời điểm thích hợp hơn.

Và trong thâm tâm, chúng tôi chắn sẽ quay trở lại.

\ct

Chúng tôi vừa bước ra khỏi bìa rừng, đồng hồ đã chỉ 10 giờ 30. Trước mặt, thị trấn Aogami hiện ra với những ánh đèn đường màu vàng ấm, quen thuộc đến nỗi tôi nghe được tiếng tim mình như vừa được tháo dây căng. Mới chỉ hơn một tiếng trước, chúng tôi còn đứng giữa khu rừng đen kịt, đối mặt với những đứa trẻ không mặt; giờ đây, mũi giày lại đạp lên mặt nhựa bằng phẳng, nghe mùi nhựa đường hòa vào gió đêm. Không khí trong nhóm nặng hẳn, như thể có bóng ma vẫn còn bám theo từng bước chân.

Mỗi người mang theo một khoảnh khắc riêng khi bước khỏi rừng.

Honoka-san là người đầu tiên cất tiếng. Cậu ấy lên tiếng thật nhẹ, như sợ đánh thức điều gì: ``\dots Tớ không nghĩ là thật sự có ma đâu.'' Giọng cậu ấy cố giữ bình thản, nhưng tôi vẫn bắt gặp ánh mắt lén liếc về phía rừng, như thể chỉ cần quay đầu lại, những bóng kia sẽ lập tức nhảy ra.

Aku-san thì còn thảm hại hơn. Cô nàng dán sát vào Shion-san, ngón tay siết chặt tay áo bạn, bước đi loạng choạng như người say. ``Chân mình vẫn còn run đây này\dots'' cô lầm bầm, giọng run rẩy. ``Lúc nãy\dots\ mấy cái bóng đó\dots\ tụi nó\dots\ tụi nó\dots'' Cô nuốt nước bọt, không nói nổi thành câu.

Shion-san vỗ vai cô bạn, hất mái tóc bạc khỏi trán, cố tỏ ra ung dung: ``Hừm. Đừng lo.'' Nụ cười nửa miệng của cô lúc này trông thật cứng đờ. ``Nếu chúng thật sự là linh hồn lang thang\dots\ thì ta sẽ kết nạp chúng làm thuộc hạ.''

Cô bạn tóc hồng-xanh giật mình, mặt cô tái nhợt, rồi hét như trời giáng: ``\textbf{Cậu đừng có nói như vậy nữa!! Tớ tưởng là chúng ta bị ma bắt đi rồi chứ!}'' Tiếng hét của cô vang lên giữa đêm khuya, khiến cả nhóm giật mình. Honoka-san bật cười, nhưng tiếng cười cũng run run, như thể chỉ cần một cơn gió thổi qua, nó sẽ vỡ tan thành mảnh.

Azuki-san vẫn quay lại nhìn, mắt còn đọng ánh đèn pin lẫn trong sương. Cô thì thào như kẻ vừa trốn thoát khỏi trang truyện tranh: ``Thật sự có thứ như thế\dots\ không phải mình tưởng tượng.'' Giọng cậu ấy rung nhẹ, nhưng không phải run sợ, mà run vì kích thích, cảm giác của đứa trẻ lần đầu thấy cá voi nhảy lên giữa biển đêm.

Touka-san bước bên cạnh, đầu cúi thấp, tay vẫn siết chặt tờ phác thảo nhàu nát. Cậu lẩm nhẩm điều gì đó về sai số chuẩn, nhưng khoé miệng giật giật: tất cả dữ liệu cô từng nghi ngờ vừa được xác nhận bằng một cái nhìn xanh lục. ``Đúng\dots\ tính toán của mình không hề lệch,'' cậu ấy thở ra, bụi sương tan trước mặt, như thể vừa tháo được nút chai chứa đầy khí đốt của riêng mình.

Saki-san lặng đến mức tiếng dép cô chạm nhựa cũng trở thành thứ duy nhất vang lên. Cô không cất lời, chỉ cúi xuống vuốt nhẹ một bông bỉ ngạn ven lề, ngón tay dừng lại ở nhụy hoa, nơi màu trắng chuyển sang vàng nhạt. Có lẽ cô đang nghĩ đến độc tố colchicin, hoặc đến việc loài hoa này vẫn nở dù không còn ai ngắm; khoảng im lặng của cô nặng như một bản báo cáo chưa viết xong.

Saya-san gấp tấm bản đồ cổ thành hình vuông nhỏ, mỗi nếp gấp đều vuông vức như cắt bằng thước. Đôi mắt cô không chớp, nhìn thẳng về phía trước, một vẻ bình thản của người đã thấy quá nhiều mảnh ghép, và biết rằng mảnh vừa rơi xuống chỉ là phần nhỏ của bức tranh chưa được phơi sáng.

Tôi bước chậm, để bốn người kia đi trước. Trong đầu tôi, đôi mắt xanh lục vẫn mở to giữa màn đêm, như ngọn hải đăng bị lạc giữa rừng. Không chỉ tôi, Onii-sama cũng thỉnh thoảng ngoái đầu, như thể muốn chắc rằng cái nhìn ấy không bám theo sau lưng. Onee-sama khoanh tay, mắt nhìn xuống đường kẻ vạch trắng, nhưng đôi chân thì bước không đều, lâu lắm rồi chị tôi mới tỏ vẻ bối rối như vậy kể từ khi được Onii-sama đi cùng ở The Void.

Sakura-san lùi về sau cùng, bước chân nhẹ đến mức không có tiếng; cậu ấy nắm chặt vòng hoa bỉ ngạn đã héo, như thể đó là chiếc chìa khóa vừa mở được cánh cửa không có then.

Chúng tôi không nói với nhau, nhưng cùng hiểu: cô bé ấy không phải là linh hồn lẫn vào đám trẻ kia. Em đứng ngoài vòng tròn, ngoài cả thời gian; một đốm xanh lục lẻ loi giữa biển đen, như người dẫn đường đã chờ từ kiếp trước, và giờ vừa gật đầu với chúng tôi: ``Các người đã đến đúng nơi rồi đấy.''

Khi ra tới ngã ba gần khu dân cư, nhóm chúng tôi bắt đầu tách ra.

Honoka-san cùng nhóm bạn thân lặng lẽ rẽ về khu dân cư phía Tây, bóng họ nhạt dần dưới đèn đường vàng vọt. Azuki-san chợt ngoái lại, vẫy tay như thể sợ đêm sẽ nuốt chữ ``Mai gặp lại nhé!'' còn vang trên môi. Cô khẽ nhắc, giọng hạ xuống thành tiếng gió: ``Nhớ giữ bí mật vụ này đó!'' lời dặn dò dù mỏng tan nhưng sức nặng của những gì đã xảy ra là quá đủ để cả nhóm gật đầu không cần tiếng.

Shion-san thậm chí còn nói nửa đùa nửa thật, mắt nheo lại như mèo đêm: ``Nếu tin này lọt ra, ngày mai cả thị trấn sẽ kéo nhau vào rừng, biến nơi ấy thành chợ.'' Aku-san nghe xong giật mình, vội vã phản đối, giọng vỡ ra thành từng mảnh: ``Xin cậu đấy, Shion\dots\ Đừng quay lại đó nữa\dots! Tớ sắp tè ra quần rồi này\dots'' Tiếng cô rung rung, nhưng cũng đủ khiến cả đám bật cười, một thứ cười giải tỏa cuối cùng trước khi đêm chính thức buông màn.

Tiếng cười vừa tắt, không gian chợt rộng ra. Chỉ còn bốn bóng người trên con đường trải nhựa về nhà: tôi, Onii-sama, Onee-sama và Sakura-san. Đèn đường cách quãng, vài cửa hàng đã kéo cửa sắt, tiếng lanh canh khuya khoắt. Gió đêm thổi nhẹ, luồn qua mái hiên, hắt hơi lạnh lên gáy, như nhắc chúng tôi: chuyện chưa kết thúc, chỉ mới tạm gác lại giữa phố Aogami yên ắng giữa màn đêm thanh tịnh.

Không ai trong chúng tôi cất lên một lời nào. Đúng hơn, không ai trong chúng tôi biết bắt đầu từ đâu. Một cuộc phiêu lưu để chứng minh lời đồn đại trong trường, hóa ra lại thành sự thật, tới mức những người yếu bóng vía còn không thể gượng nổi.

Đúng lức ấy, giữa phố đêm tĩnh mịch, một âm thanh trầm ấm vang lên từ cổ tay Onii-sama: Game-san, giọng Pháp khản đặc đủ khiến cả nhóm đứng sững.

``Tôi đã xác định được số lượng linh hồn,'' ông ta nói.

Chúng tôi dừng bước.

Game tiếp tục: ``Nhưng giải phóng chúng\dots\ sẽ không hề đơn giản.''

Onee-sama nhíu mày. ``Bao nhiêu?''

``Có ít nhất mười hai linh hồn đang mắc kẹt.'' Con số làm tôi thót tim; tôi tưởng chỉ vài bóng trẻ co thôi, ai ngờ đông gấp bội. Ông ta gầm gừ: ``Tha cho họ không dễ; cần một nghi thức đúng chuẩn và\dots\ một lời xin lỗi muộn màng.''

Ông ta ``hừm'' một tiếng, nghe như thở dài, rồi quay lại hỏi chúng tôi: ``Có ai trong các cậu nghĩ ra được kế nào không?''

Không ai đáp. Onii-sama ngước nhìn trời, mắt ánh lên vẻ ``lại thêm bài toán cổ''. Onee-sama khoanh tay, môi mím chặt như đang giải mã khối rubik bằng ý nghĩ. Sakura-san quay về phía rừng, tay siết bông bỉ ngạn héo, có lẽ cậu ấy đang nghe tiếng hoa thì thầm. Tôi thì chìm trong đôi mắt xanh lục kia: cô bé đứng ngoài thời gian, vừa gật đầu với tôi, vừa lặng im như hồ nước đóng băng.

Cuối cùng, Onii-sama phá vỡ im lặng: ``Ta phải hiểu nghi thức xưa ấy đã vận hành thế nào, và vì sao lũ trẻ chưa thể được an nghỉ.''

Onee-sama gật nhẹ: ``Có lẽ lỗi không phải của chúng, mà của kẻ đã khép cửa.'' Sakura-san thì thào: ``Xem chừng nói chuyện với chúng nó\dots\ là cách đầu tiên mà tôi có thể nghĩ tới. Ý tôi là, chẳng phải hỏi trực tiếp chúng sẽ dễ hơn so với việc phải lục soát thông tin sao?''

``\dots Cậu nói cũng phải,'' tôi nói. Chắc gì thư viện trường chúng ta có ghi chép gì về lễ tế này, mà cả thành phố - theo lời của Akane-san - cũng không có lấy một thư viện lớn nào hay bảo tàng nào cả, huống gì là một nơi có thể lưu trữ.

Game-san im lặng ba giây, rồi cất tiếng như khép lại một bản nhạc lỗi: ``Vậy thì ta còn việc phải làm đấy. Không ai biết được điều gì sẽ xảy đến với chúng ta trong những ngày tiếp theo đâu.''

Chúng tôi bước đi, nhịp bước vang trên mặt nhựa lạnh. Phố Aogami mở ra trước mắt như một dải lụa vàng nhạt, đèn đường tỏa quầng sáng tròn trĩnh, gió đêm luồn qua mái ngói im lìm, rồi lướt nhẹ lên tóc chúng tôi như kẻ thở dài.

Một tuần rồi, đúng bảy ngày kể từ khi chúng tôi bắt đầu một cuộc sống ở thị trấn Aogami mới. Không còn tiếng kim đồng hồ giật lùi, không còn ánh mắt quen thuộc nào lặp đi lặp lại trong gương. Mặt trời vẫn mọc đúng giờ, phút vẫn trôi đều đặn, thời gian như dòng nước mùa xuân ấm lên, cuốn theo bụi bẩn của ngày hôm qua.

Có lẽ ta nên mừng: cuộc sống đã trở lại quỹ đạo vốn có ở vùng đất này. Nhưng mỗi tối, khi đèn phòng bật sáng, tôi vẫn thấy bốn bóng người ngồi thành vòng, im lặng như đang nghe tiếng lá rơi trong đầu. Những linh hồn khu rừng chưa buông chúng tôi; họ đi cùng về nhà, len lỏi trong giấc ngủ, gõ nhẹ trán mỗi khi đồng hồ đã điểm tới nửa đêm.

Và đôi mắt xanh lục kia - trong veo, lạnh lẽo, đậu lại trên cửa sổ như con bướm đêm - vẫn mở to, chờ một lời hứa chưa thành tiếng.

Trước mắt, tôi có thể thấy em ấy - và những linh hồn trẻ con ấy - không có ác tính gì. Nhưng tôi không dám chắc. Phải chăng những linh hồn đó\dots\ bị mắc kẹt lại? Liệu chúng còn những điều gì đó chưa được thỏa mãn không?

\pov{-}

Bốn bóng người vừa khuất khỏi lối mòn, khoảng đất trước đền liền chìm vào im lặng như hồ nước tạnh sau cơn mưa. Đèn lồng giấy đung đưa, tỏa vòng sáng loang lổ trên bậc đá rêu phong. Một làn gió thổi qua, mang theo hơi lạnh của lá ướt và mùi đất sâu, khiến ngọn đèn run rẩy như vừa tỉnh giấc sau giấc ngủ dài.

Từ lớp sương mỏng tan dần, một thân hình nhỏ xuất hiện. Cô bé khoác kimono nâu sẫm thời Nara, tà áo thêu hoa cúc đã bạc màu; tóc đen tết đuôi sam khẽ bay, phơi những sợi chỉ âm u dưới ánh đèn mờ. Đôi mắt xanh lục như hồ nước ngọt giữa rừng sâu, lặng lẽ hướng về con đường vừa nuốt chửng bốn kẻ lữ khách.

Không buồn, không giận, chỉ là một vẻ trống rỗng mỏng manh. Cô bé đứng yên trên bậc thềm đá, lắng nghe tiếng lá rơi, tiếng gió rít, tiếng trái tim mình đập nhịp xa xăm. Trong đáy mắt ấy, bốn bóng hình lạ lẫm vẫn còn in hằn: cậu con trai tóc rối như tổ quạ, đôi mắt đỏ sâu thẳm như đêm cuối cùng của năm; cô gái tóc đuôi ngựa xù, bước đi như luôn sẵn sàng đi ngược lại cả thế giới; cô gái mắt hai màu xanh-đỏ, nhìn mọi thứ như thể đang đọc một bản nhạc chỉ mình cô nghe được; và cô nàng tóc hồng, tay siết bông hoa trắng như giữ lấy mảnh ký ức cuối cùng của mình. Bốn người ấy đã đi rồi, nhưng dư âm của họ vẫn còn vương trong không gian, như hương hoa bỉ ngạn thoang thoảng không tan.

Cô bé khẽ nghiêng đầu, mái tóc đen trượt xuống vai, che đi nửa khuôn mặt nhỏ nhắn. Đôi mắt xanh lục như mở ra một cửa sổ khác, nơi thời gian không còn là đường thẳng mà là vòng tròn khép kín. Bốn kẻ kia không biết tên cô, cũng như cô không biết tên họ; nhưng trong khoảnh khắc họ đứng trước mặt, cô đã cảm nhận được điều gì đó mà bao năm qua chưa từng có: một rung động nhỏ, như tiếng đàn xa vọng, như hơi ấm từ bàn tay vô tình chạm phải làn da ma thuật của cô.

Đôi môi nhỏ mấp máy, giọng thì thầm tan trong không khí, không phải nói với ai, mà như nói với chính mình, với khu rừng, với thời gian đã ngừng đập từ lâu:

\begin{center}
  ``Cuối cùng\dots\ cũng có người nghe thấy chúng ta rồi.''
\end{center}

Câu nói vừa dứt, gió chợt thổi mạnh hơn, đèn lồng lắc lư như muốn rời khỏi móc treo. Khi ánh sáng định hình trở lại, bậc đá chỉ còn bóng đêm phủ kín. Cô bé đã biến mất, để lại hương gỗ cổ kính thoang thoảng và tiếng lá xào xạc như một bản nhạc nền không lời, vang lên mãi trong khu rừng tĩnh mịch.

\end{document}